\subsection{A compact limit and its global balances}

Taking $q=2-a$ in \eqref{eq:entropy-global} makes
$\Psi_{k^q}$ a positive multiple of $k$; taking $q=2$ controls the
viscous gradient. The boundary errors are $O(\eps_j^\beta)$ and the
outer weights are uniformly bounded. Thus
\begin{equation}\label{eq:compactness-bounds}
 \sup_j\int_0^T\|w_j\|_{H^{\alpha/2}(\R^3)}^2\,\dd t<\infty,
 \qquad \sup_j\eps_j\int_0^T\int_{D_j(t)}|\nabla w_j|^2<\infty.
\end{equation}
The pressure flux is bounded on a common enclosing ball. The zero-extended
viscous flux tends to zero in $L^2$, and \eqref{eq:extended-mass} gives
\begin{equation}\label{eq:time-bound}
 \partial_tw_j\ \hbox{bounded in }L^2_tH^{-1}_x,
 \qquad \|w_j(t)-g_j\|_{H^{-1}}\leq K(t+\sqrt{\eps_jt}).
\end{equation}
Here $H^{-1}$ is the dual of $H^1_0$ on a slightly larger fixed ball;
no whole-space $H^1$ norm across the outer jump is used.

The compactness lemma \cite[Lemma~7.5]{Stan2016} applies with
$X=H^{\alpha/2}$, $B=L^2$, $Y=H^{-1}$ and exponent two. Rellich gives
the first compact embedding, and Poincar\'e gives the second continuous
one. The derivative bound implies the required $L^1$-in-time bound on
finite intervals. A subsequence therefore satisfies
\begin{equation}\label{eq:limit-convergence}
 w_j\to u\quad\hbox{strongly in }L^2([0,T]\times\R^3),
 \qquad R_j\to R\quad\hbox{uniformly}.
\end{equation}
For smooth spatial tests, \eqref{eq:time-bound} gives equicontinuity.
A diagonal extraction on a countable dense family selects a density
$u(t)$ at every time, with $0\leq u(t)\leq U$. Approximate any bounded
spatial test in $L^1$ on the common support by a smooth test; the cap
bounds both approximation errors uniformly in time. Hence
\begin{equation}\label{eq:weak-time}
 \sup_{0\leq t\leq T}\left|\int(w_j(t)-u(t))\psi\right|\to0
 \quad(\psi\in L^\infty),\qquad
 u(0)=F,\quad \int u(t)=\int F.
\end{equation}
All such pairings are continuous. The uniformly convergent radii imply
$u(t)=0$ outside $B_{R(t)}$ at every time.

Truncate the near singularities of the potential and velocity. The retained
kernels on the common support form a compact family of spatial $L^1$
tests as the evaluation point varies on a compact set. A finite cover,
\eqref{eq:weak-time} and the near-field bounds give
\begin{equation}\label{eq:field-limit}
 K_\alpha w_j\to K_\alpha u,\qquad v[w_j]\to v[u]
 \quad\hbox{uniformly on compact space-time sets}.
\end{equation}
The Stefan law and \eqref{eq:boundary-flux} then yield
\begin{equation}\label{eq:radius-limit}
 R_j\to R\quad\hbox{in }C^1([0,T]),\qquad
 R'=A^{a-1}v[u](R),\qquad R(0)=R_0.
\end{equation}
The layer limit gives $u\geq c(t)\geq A/2$ inside its ball and
\begin{equation}\label{eq:trace-modulus}
 |u(t,r)-A|\leq\frac{K}{\kappa\beta}(R(t)-r)^\beta
 \quad\hbox{near }R(t).
\end{equation}
These inequalities pass by bounded-test pairings, so hold at every time
in the essential spatial sense. In particular the inner trace is $A$.
Strong density convergence and \eqref{eq:field-limit} pass the nonlinear
flux in local $L^1$; the viscous flux vanishes. The same $u$ solves the
whole-space equation with initial $F$ and locally integrable flux.

The energies in \eqref{eq:pressure-identity} converge at every terminal
time by \eqref{eq:weak-time}, \eqref{eq:field-limit} and
$\mathcal E_\alpha(w)=\tfrac12\int wK_\alpha w$.
The pressure dissipation converges in $L^1$ on time intervals by bounded
fields and strong density convergence. Its positive error
$\eps_j A\int_0^tW[w_j;R_j]\dd r$ vanishes, while the other viscous term
is nonnegative. Therefore
\begin{equation}\label{eq:limit-pressure}
 \mathcal E_\alpha(u(t))+
 \int_0^t\int u^a|\nabla K_\alpha u|^2
 \leq\mathcal E_\alpha(F),\qquad 0\leq t\leq T.
\end{equation}

To retain the viscosity in the power balances, extract once the joint
nonnegative finite measures
\begin{equation}\label{eq:nu}
 \nu_j=\eps_j|\nabla w_j|^2\mathbf1_{D_j(t)}\,\dd x\dd t\,
             \delta_{w_j(t,x)}(\dd k)\ \rightharpoonup\ \nu.
\end{equation}
Their mass is bounded by \eqref{eq:compactness-bounds}; they have common
compact space-time support and state support $[A/2,U]$.
Weak compactness follows, for example, by diagonal extraction on a
countable dense family of continuous functions on that compact set and
extension of the resulting positive bounded functional. This extraction
is independent of the entropy exponent.

At good terminal times, \eqref{eq:entropy-global} passes to
\begin{align}
 &\|u(t)\|_q^q+\int_0^t\mathcal B_\alpha(\Psi_{k^q},u(r))\,\dd r\nonumber\\
 &\quad+q(q-1)\int_{(0,t)\times\R^3\times[A/2,U]}k^{q-2}\,\dd\nu\nonumber\\
 &\leq\|F\|_q^q+q(q-1)\int_0^tW[u(r);R(r)]\,\dd r
             \int_0^A k^{q-2}g_A(k)\,\dd k.\label{eq:limit-power-cost}
\end{align}
A good time has strong spatial convergence along the subsequence and no
time atom of $\nu$. Fatou gives the positive fractional term, measure
convergence gives the state moment, and \eqref{eq:field-limit} passes the
surface weight. The boundary error vanishes. Take first a countable dense
set of $q>1$ and intersect the good-time sets. Extension to real $q$ uses
bounded endpoint powers, continuity of $k^{q-2}$ on $[A/2,U]$ and Fatou
for the positive pair integrands. Near zero, the outer density integrals
are dominated by a multiple of
$k^{a+q_{\min}-2}+k^{q_{\min}-1}$ in a neighborhood of each $q>1$.

In this passage initially include time zero in the measure moment. The
quadratic case gives, at good times,
\[
 \|u(t)\|_2^2\leq\|F\|_2^2+Kt,\qquad
 2\nu([0,t]\times\R^3\times[A/2,U])
 \leq\|F\|_2^2-\|u(t)\|_2^2+Kt.
\]
Weak-time lower semicontinuity extends the first bound to every time.
The initial weak trace and cap imply weak $L^2$ convergence to $F$ on
the common support. The norm bound makes it strong as $t\downarrow0$.
The second bound now excludes an initial atom of $\nu$, justifying
$(0,t)$ in \eqref{eq:limit-power-cost}.

There is a fixed annulus about $S$ on which the velocity remains positive
for a common interval beginning at zero, uniformly in the retained
approximants. Indeed \eqref{eq:time-bound} controls smooth retained
velocity kernels; their omitted near kernels are uniformly small. A finite
test cover gives a common time modulus near $S$, and the initial fields
converge uniformly there to $v[F]$. Its value at $S$ is positive and its
radial component is continuous. Shorten $T$ and discard a finite prefix
so that an annulus $J_+$, separated from zero and all outer boundaries,
satisfies
\begin{equation}\label{eq:positive-annulus}
 v[w_j](t,r),\ v[u](t,r)\geq v_*>0,
 \qquad (t,r)\in[0,T]\times J_+.
\end{equation}
Finite deletion changes neither limit. A restriction time with no terminal
atom of $\nu$ can be chosen. We use this same restricted density, radius,
tail sequence and measure in the rest of the proof.
